% 19 November 2026. Drafted for Edward Baker with OpenAI Codex assistance.
% Helper explicitly dispatched as gpt-6.1-sol, reasoning effort ultra;
% independent serving-backend attestation was not available to the helper.
\section{The full arithmetic reference and its useful domain}
\label{sec:arithmetic-source}

The arithmetic input to a backward continuation need not be reduced to
the two leading carriers that produce an asymptotic clock.  At a positive
time the finite Riemann--Siegel approximation retains many arithmetic
terms, with their phases, before its error is estimated.  The question
for this program is whether those terms can supply effective gap data
where the two-carrier argument no longer does so.  This is a question
about an additional quantitative estimate.  The approximation and the
identities below do not themselves extend the known all-real region.
We use the published approximation of \cite{Polymath2019} and the
selected-source analysis in \cite{VCDCMDiscussion11,VCDCMDiscussion12}.

\subsection{A holomorphic finite reference}

To avoid confusing an arithmetic carrier with the complete exterior
interaction of a zero pair, write $\mathfrak M_t$ for the function called
$M_t$ in \cite{Polymath2019}.  On
$\mathcal D=\C\setminus(-\infty,1]$, use the principal logarithms to set
\begin{align}
 L_0(s)&=\operatorname{Log}s+\operatorname{Log}(s-1)
       -\frac{s}{2}\log\pi+\log\frac{\sqrt{2\pi}}{16}
       +\left(\frac{s}{2}-\frac12\right)
           \operatorname{Log}\frac{s}{2}-\frac{s}{2},
       \label{as:carrier-log}\\
 \mathfrak M_0(s)&=\exp L_0(s),\qquad
 \alpha(s)=L_0'(s)
       =\frac{1}{2s}+\frac{1}{s-1}
         +\frac12\operatorname{Log}\frac{s}{2\pi},
       \label{as:alpha}\\
 \mathfrak M_t(s)&=\mathfrak M_0(s)
                  \exp\left(\frac{t}{4}\alpha(s)^2\right).
       \label{as:carrier}
\end{align}
These carriers are nonvanishing on $\mathcal D$.  The logarithm in
\eqref{as:carrier-log} fixes an actual analytic function, rather than a
pointwise choice of argument.  For $n\in\N$ define
\begin{align}
 \ell_n(s)&=L_0(s)-s\log n,\qquad
 a_n(s)=\ell_n'(s)=\alpha(s)-\log n,\\
 p_n(t,s)&=\exp\left(\ell_n(s)+\frac{t}{4}a_n(s)^2\right).
       \label{as:pn}
\end{align}
Expanding the square gives the useful exact identity
\begin{equation}
 p_n(t,s)=\mathfrak M_t(s)\,
     \exp\left(\frac{t}{4}\log^2 n\right)
     n^{-s-(t/2)\alpha(s)}.
 \label{as:pn-dirichlet}
\end{equation}
In particular, the $n=1$ term is $\mathfrak M_t(s)$, and all the terms
with $2\le n\le N$ remain part of the reference.

Put
\begin{equation}
 s_+(z)=\frac{1+\ii z}{2},\qquad
 s_-(z)=\frac{1-\ii z}{2},\qquad
 P_N(t,z)=\sum_{n=1}^N\{p_n(t,s_+(z))+p_n(t,s_-(z))\}.
 \label{as:PN}
\end{equation}
For a fixed integer $N$, this is holomorphic in $z$ on the right
half-plane.  Indeed $\Re z>0$ places the two $s$-arguments in opposite
open half-planes, away from the cut of $L_0$.  For real $t$ it satisfies
$P_N(t,\overline z)=\overline{P_N(t,z)}$ there.  It therefore has the
local conjugation symmetry required of a reference for a real zero pair.
This assertion about its definition is separate from the domain of a
published error bound.

For $z=x+\ii y$ with $x\ge200$ and $0\le y\le1$, the convenient
nonvanishing normalizer is
\begin{equation}
 \mathcal B_t(z)=\mathfrak M_t(s_-(z)),\qquad
 \gamma_t(z)=\frac{\mathfrak M_t(s_+(z))}
                  {\mathfrak M_t(s_-(z))},\qquad
 b_n^t=\exp\left(\frac{t}{4}\log^2 n\right).
 \label{as:normalizer}
\end{equation}
Consequently the normalized full reference is
\begin{equation}
 f_{t,N}(z)=\frac{P_N(t,z)}{\mathcal B_t(z)}
   =\sum_{n=1}^N\frac{b_n^t}{n^{s_-(z)+(t/2)\alpha(s_-(z))}}
    +\gamma_t(z)\sum_{n=1}^N
       \frac{b_n^t}{n^{s_+(z)+(t/2)\alpha(s_+(z))}}.
 \label{as:full-f}
\end{equation}
This form makes holomorphy visible.  A formula written using $x,y$ and
complex conjugates must first be rewritten in this form before complex
spatial differentiation is justified; see \cite[\S8, (92)]{Polymath2019}.

\subsection{The exact arithmetic terms behind the approximation}

The finite reference comes from an exact expansion, rather than an
independent model of the zero divisor.  Define
\begin{equation}
 K(s)=\frac{s(s-1)}{16}\pi^{-s/2}\Gamma(s/2),\qquad
 r_{0,n}(s)=K(s)n^{-s}.
 \label{as:exact-initial-term}
\end{equation}
For nonreal $s$ and $t>0$ the Gaussian evolution of this term is
\begin{equation}
 r_n(t,s)=\frac1{\sqrt\pi}\int_{\R}
          r_{0,n}(s+\sqrt t\,v)e^{-v^2}\,\dd v.
 \label{as:exact-term}
\end{equation}
The horizontal integration line avoids the poles of the meromorphic
factor $K$.  On compact sets separated from the real $s$-axis, its
growth along that line is slower than a Gaussian, so the integral and
the differentiations used below converge locally uniformly.

For completeness, the exact Riemann--Siegel remainder before Gaussian
evolution is
\begin{equation}
 R_{0,N}(s)=K(s)\int_{N\swarrow N+1}
       \frac{w^{-s}e^{\pi\ii w^2}}
            {e^{\pi\ii w}-e^{-\pi\ii w}}\,\dd w.
 \label{as:exact-RS-remainder}
\end{equation}
Here the oriented diagonal line has direction $e^{5\pi\ii/4}$ and
crosses the real axis between $N$ and $N+1$; the conventional logarithm
along the Riemann--Siegel contour is understood.  Let $R_N(t,s)$ be its
Gaussian evolution as in \eqref{as:exact-term}, and write
$F^*(s)=\overline{F(\overline s)}$.  The exact decomposition is
\begin{align}
 H_t(z)={}&\sum_{n=1}^N r_n(t,s_+(z))
        +\sum_{n=1}^N r_n^*(t,s_-(z))\notag\\
        &+R_N(t,s_+(z))+R_N^*(t,s_-(z)).
 \label{as:exact-finite-expansion}
\end{align}
It holds for every fixed nonnegative integer $N$, $t>0$, and
$\Re z\ne0$; this is the Gaussian form of the Riemann--Siegel identity
in \cite[\S4]{Polymath2019}.  The star is retained to display the
reflection in the exact formula.  The coefficients and the Gaussian
kernel are real under conjugation, so that $r_n^*=r_n$ on the two
nonreal components.  Replacing each exact $r_n$ by \eqref{as:pn} gives
\eqref{as:PN}; neither replacement removes the complete remainder.

The approximation of an individual term has an explicit, relative
source bound.  With $s=\sigma+\ii T$, $T>10$, and $0<t\le1/2$, one may
write
\begin{equation}
 r_n(t,s)=p_n(t,s)(1+\delta_n(t,s)),\qquad
 |\delta_n(t,s)|\le
 \exp\left(
 \frac{(t^2/8)|\alpha(s)-\log n|^2+t/4+1/6}{T-3.33}
 \right)-1.
 \label{as:term-source-bound}
\end{equation}
The corresponding lower-half-plane statement follows by conjugation.
This is the term estimate of \cite[Proposition 6.1]{Polymath2019}; its
error is relative to an individual carrier term, rather than to $H_t$
near one of its zeros.  Summing it and paying the Riemann--Siegel
remainder supplies the complete approximation error.

The standard cutoff in that complete bound is
\begin{equation}
 N(t,x)=\left\lfloor\sqrt{\frac{x}{4\pi}+\frac{t}{16}}\right\rfloor.
 \label{as:natural-cutoff}
\end{equation}
Theorem 1.3 of \cite{Polymath2019} gives, throughout
\begin{equation}
 0<t\le\tfrac12,\qquad x\ge200,\qquad 0\le y\le1,
 \label{as:source-domain}
\end{equation}
the complete bound
\begin{equation}
 H_t(z)=P_{N(t,x)}(t,z)+E_0(t,z),\qquad
 |E_0(t,z)|\le |\mathcal B_t(z)|
             \{e_A(t,z)+e_B(t,z)+e_{C,0}(t,z)\}.
 \label{as:source-value-bound}
\end{equation}
The three $e$'s in this display denote the explicit positive budgets of
that theorem, with its definitions (71)--(74); they are a named imported
input, not quantities bounded in this manuscript.  The first two pay
the term replacements and the last pays the whole remainder when its
leading correction is omitted.  No numerical conclusion below requires
an unevaluated choice of these budgets.  An application must use their
actual formulas, with the same normalization and domain.  Bounds below
the real axis can be supplied by conjugating the physical approximation.

Two qualifications matter when using \eqref{as:source-value-bound} in
a differential argument.  First, the theorem does not impose
$x\ge\exp(C/t)$.  Its domain is larger than the clock regime.  Second,
the natural cutoff jumps.  One must use a time and spatial cell on
which \eqref{as:natural-cutoff} has one value, or keep a fixed $N$ and
explicitly transfer the omitted or added terms to the error.  For example,
at the same $(t,z)$,
\begin{equation}
 P_{N+1}(t,z)-P_N(t,z)
       =p_{N+1}(t,s_+(z))+p_{N+1}(t,s_-(z)).
 \label{as:cutoff-jump}
\end{equation}
This is a known change of reference.  It must be included in a seam
payment if a cell crosses a cutoff surface.  Differentiating the floor
in \eqref{as:natural-cutoff} as if it were constant would instead lose
that payment.

\subsection{What the two-carrier clock discards}

On the real axis put $L=\log(x/(4\pi))$ and $m=\log n$.
Directly from \eqref{as:alpha},
$\Re\alpha((1-\ii x)/2)=L/2+O(x^{-2})$.  Therefore the magnitude of
an individual normalized term in the first sum of \eqref{as:full-f} is
\begin{equation}
 \left|b_n^t n^{-s_--(t/2)\alpha(s_-)}\right|
   =n^{-1/2}
      \exp\left\{-\frac{t}{4}m(L-m)+O(tm/x^2)\right\}.
 \label{as:clock-term-decay}
\end{equation}
For $n\le N(t,x)$, one has
$m\le L/2+O(t/x)$.  When $tL$ is sufficiently large, summing
\eqref{as:clock-term-decay} over $n\ge2$ gives a negative power of
$x$ with exponent proportional to $t$, after decreasing the absolute
constant in that exponent.  The same is true for the reflected sum.
The essential large parameter for this dominance is $t\log x$.
Accordingly a uniform small-tail argument has an onset of the form
$x\ge\exp(C/t)$, with an absolute constant that has to be large enough
for all the other estimates as well.

The coefficient $t/4$ in \eqref{as:clock-term-decay} follows from the
definitions by expanding a square.  It is not replaced here by the
$t/2$ coefficient appearing in an intermediate display of the source's
\S9; the discrepancy and its source check are recorded in
\cite{VCDCMDiscussion11}.  We extract no numerical onset constant from
this qualitative summation argument.

In the dominance region the reference reduces to
\begin{equation}
 P_N(t,z)\simeq\mathfrak M_t(s_+(z))+
                      \mathfrak M_t(s_-(z)).
 \label{as:two-carriers}
\end{equation}
Cancellation of these two carriers yields the clock
\begin{equation}
 G(x,t)=\frac{x}{4\pi}\log\frac{x}{4\pi}
       -\frac{x}{4\pi}+\frac{11}{8}
       +\frac{t}{16}\log\frac{x}{4\pi},\qquad
 G(q_n(t),t)=n.
 \label{as:clock}
\end{equation}
Theorem 1.5 of \cite{Polymath2019} supplies actual real simple zeros
near these clock points above its height and index onsets.  The existing
analysis \cite{VCDCMManuscript} uses one buffered onset on a slab
$0<t_0\le t\le t_*\le1/2$, and hence one set of common labels.
The distinction between a height onset and an index onset is preserved;
both must be met before those labels can be used.

The full reference \eqref{as:PN} remains available when the $n\ge2$
terms cannot be treated as a vanishing tail.  They can then change the
phase, location, and separation of its zeros.  They belong in the
reference field being followed.  Treating them as a new small error
without a bound would simply reintroduce the clock hypothesis.  On the
other hand, retaining them does not prove that their own divisor is
real and simple, or that the source error is smaller than a gap margin.
The useful initial datum is an enclosure of an \emph{actual} gap
$g=b-a$, or of $D=g^2/4$, with isolation and completeness paid.  A mean
spacing $4\pi/L$ becomes such a datum only after localization errors
are smaller than the relevant separation.  This is why a source formula
with a wider domain is an opportunity for continuation, rather than an
already acquired continuation theorem.

\subsection{The strip certificate and the existing upper bounds}

A global upper bound for $\Lambda$ has another consumer besides a list
of separated real pairs.  We recall the exact domains of the criterion
being compared.  In the nontrivial regime $0<t_0<1/2$, $X>0$ and
$0<y_0<\sqrt{1-2t_0}$, the three inputs of
\cite[Theorem 1.2]{Polymath2019} are the following zero exclusions:
\begin{align}
 &\zeta(\sigma+\ii T)\ne0
 &&\left(\frac{1+y_0}{2}\le\sigma\le1,
                       \quad0\le T\le X/2\right),
                       \label{as:strip-initial}\\
 &H_{t_0}(x+\ii y)\ne0
 &&\left(x\ge X+\sqrt{1-y_0^2},
                   \quad y_0\le y\le\sqrt{1-2t_0}\right),
                       \label{as:strip-final}\\
 &H_t(x+\ii y)\ne0
 &&\left\{\begin{aligned}
       &X\le x\le X+\sqrt{1-y_0^2},\quad0\le t\le t_0,\\
       &\sqrt{y_0^2+2(t_0-t)}\le y\le\sqrt{1-2t}.
       \end{aligned}\right.
                       \label{as:strip-barrier}
\end{align}
Together they give
\begin{equation}
 \Lambda\le t_0+\frac{y_0^2}{2}.
 \label{as:strip-output}
\end{equation}
These parameter restrictions keep the displayed regions meaningful
and include the applications discussed below. The last two inputs
are nonvanishing problems for the complete source, so finite arithmetic
lower bounds and their full error budgets can be useful without resolving
every real gap.

The extra $y_0^2/2$ is the payment for contracting a global zero strip
to the real axis.  In particular, \eqref{as:strip-final} and
\eqref{as:strip-barrier} do not assert that all zeros are already real
at $t_0$.  A real-pair backward argument beginning at $t_0$ cannot use
that strip statement as an all-real, separated initial divisor.  A later
all-real slice and the quantitative separation needed by its chosen
comparison are additional data.  Even a global all-real statement by
itself is not a numerical lower bound for every gap.

Platt and Trudgian prove finite RH verification through
$T_{\rm PT}=3\,000\,175\,332\,800$.  Their Corollary 2 obtains
$\Lambda\le0.2$ by applying the existing Polymath parameter table to
that verified height \cite{PlattTrudgian2021}.  The conversion is $X=2T$.
The result supplies complete finite zero exclusion at time zero, rather
than the fine individual gap enclosures proposed here.

At the pinned revision of Gomila's proof package, the reported row is
\begin{equation}
 X=6\,000\,000\,185\,827,\qquad t_0=\frac{129}{800},\qquad
 y_0^2=\frac{87677}{2500000}.
 \label{as:gomila-row}
\end{equation}
Its claimed computer-assisted unconditional result, awaiting peer
review, is $\Lambda\le0.1787854$ \cite{Gomila2026Proof}.  The exact
parameter arithmetic is
\begin{equation}
 t_0+\frac{y_0^2}{2}=0.16125+0.0175354=0.1787854,
 \qquad
 T_{\rm PT}-\frac X2=\frac{350479773}{2}>0.
 \label{as:gomila-arithmetic}
\end{equation}
Thus the exact verified endpoint covers the requested initial height.
Using the rounded headline $3\cdot10^{12}$ would obscure that reserve.
This check is parameter bookkeeping, not an independent replay of the
proof package.  Its finite arithmetic lower bounds, infinite-tail
argument, and barrier over the intervening time interval are separate
certificate obligations.  The repository audit has a stated partial
replay scope \cite{NewmanCertificateAudit}; neither that audit nor the
present exposition upgrades the external package's review status.

The live explanatory page \cite{Gomila2026Blog} also reports a later
0.1782354 calculation at the same verified height, with its mathematical
review still pending. That separate research report is not a bound
verified by this investigation and is not substituted for the pinned
package compared here.

These comparisons specify what a future gain must accomplish.  A local
improvement in a gap continuation margin can be useful, but it is not
yet an improvement of \eqref{as:strip-output}.  The latter also needs
complete height coverage and the correct initial, final, and barrier
interfaces.  Conversely a better strip certificate need not supply the
separated all-real starting data of a gap program.  The two objectives
share arithmetic source information and pay different final margins.
